% Planificación Real · Los Gantt-sters · MDA/DGP 2026-2027
\documentclass[11pt,a4paper]{article}

\usepackage{iftex}
\ifPDFTeX
 \usepackage[T1]{fontenc}
 \usepackage[utf8]{inputenc}
 \usepackage{tgheros}
 \renewcommand{\familydefault}{\sfdefault}
\else
 \usepackage{fontspec}
 \setmainfont{TeX Gyre Heros}
 \setmonofont{Latin Modern Mono}
\fi
\usepackage[spanish,es-noshorthands,es-tabla]{babel}
\usepackage[a4paper,top=2.9cm,bottom=2.4cm,left=2.5cm,right=2.5cm,headheight=20pt,headsep=0.8cm]{geometry}
\usepackage[table]{xcolor}
\usepackage{needspace,graphicx,longtable,array,enumitem,titlesec,fancyhdr,caption}
\usepackage{etoolbox}
\usepackage{amsmath}
\usepackage[most]{tcolorbox}
\usepackage[hidelinks]{hyperref}

% ---------- Identidad visual común ----------
\definecolor{gblue}{HTML}{1F4FB4}
\definecolor{gnavy}{HTML}{1B2236}
\definecolor{ggreen}{HTML}{2E9E5B}
\definecolor{ggray}{HTML}{5F6B7A}
\definecolor{gpale}{HTML}{8D919B}
\definecolor{glight}{HTML}{E9EFFB}
\definecolor{gline}{HTML}{B1C1E5}

\newcommand{\proyecto}{Planificación Real · MDA}
\newcommand{\hl}[1]{\textcolor{gblue}{\textbf{#1}}}

\color{gnavy}
\setlength{\parindent}{0pt}
\setlength{\parskip}{6pt}
\renewcommand{\baselinestretch}{1.1}
\hypersetup{pdftitle={Planificación Real – Los Gantt-sters},pdfauthor={Los Gantt-sters},pdfsubject={Planificación Real y Seguimiento del Proyecto}}

% Títulos de sección: número en recuadro azul, texto azul y filete verde
\titleformat{\section}{\Large\bfseries\color{gblue}}
 {\colorbox{gblue}{\textcolor{white}{\thesection}}}{0.6em}{}[\vspace{-4pt}{\color{ggreen}\titlerule[0.8pt]}]
\titlespacing*{\section}{0pt}{18pt}{10pt}
\titleformat{\subsection}{\large\bfseries\color{ggreen}}{\thesubsection}{0.6em}{}
\titlespacing*{\subsection}{0pt}{10pt}{4pt}

\setlist[itemize]{label={\color{gblue}\textbullet},leftmargin=1.4em,itemsep=1pt,parsep=0pt,topsep=2pt}
\captionsetup{font={small,it},labelfont={small,it},textfont={color=ggray},labelsep=period}
\setlength{\LTpre}{6pt}\setlength{\LTpost}{8pt}
\renewcommand{\arraystretch}{1.35}
\arrayrulecolor{gline}

% ---------- Cajas ----------
\newtcolorbox{verde}[1]{enhanced,colback=white,colframe=ggreen,boxrule=0.8pt,arc=2pt,
 left=6pt,right=6pt,top=4pt,bottom=4pt,fonttitle=\small\bfseries\color{white},title={#1},before upper=\small}

% Cabecera y pie
\pagestyle{fancy}\fancyhf{}
\renewcommand{\headrulewidth}{0pt}
\fancyhead[L]{\raisebox{-0.25\height}{\includegraphics[height=0.55cm]{logo.png}}\hspace{0.5em}{\small\bfseries\color{gblue}LOS GANTT-STERS}}
\fancyhead[R]{\small\color{ggray}\proyecto\ ·\ Curso 2026-2027}
\fancyfoot[C]{{\color{gpale}\rule{2.5cm}{0.4pt}}\\[2pt]\small\color{ggray}\thepage}

\begin{document}
\sloppy

% ---------- Portada ----------
\begin{titlepage}
\centering
\vspace*{0.2cm}
{\fontsize{30}{34}\selectfont\bfseries\color{gblue}MDA\par}
\vspace{4pt}
{\small\bfseries\color{ggreen}Metodologías de Desarrollo Ágil\par}
\vspace{0.8cm}
\includegraphics[width=4.2cm]{logo.png}\par
\vspace{0.9cm}
{\fontsize{24}{28}\selectfont\bfseries\color{gnavy}LOS GANTT-STERS\par}
\vspace{8pt}
{\color{ggreen}\rule{\linewidth}{1.2pt}\par}
\vspace{0.8cm}
{\fontsize{20}{24}\selectfont\bfseries\color{gblue}PLANIFICACIÓN REAL\par}
\vspace{6pt}
{\bfseries\color{gnavy}SEGUIMIENTO DE HITOS Y ESTADO DEL PROYECTO\par}
\vspace{1.1cm}
{\footnotesize\bfseries\color{ggray}Integrantes\par}
\vspace{4pt}
{\color{gblue}\rule{\linewidth}{1.2pt}\par}
\vspace{10pt}
\begin{tabular}{c}
Javier Miñán Molina\\[3pt]
Germán Delgado Fernández\\[3pt]
Ana Meca Cano\\[3pt]
Jaime Montávez Ruiz\\[3pt]
Álvaro Sanabria Díaz\\[3pt]
Juan Luis Jiménez Urbano
\end{tabular}\par
\vspace{14pt}
{\small\color{gpale}Universidad de Granada · Curso 2026-2027\par}
\end{titlepage}

{\hypersetup{linkcolor=gnavy}\renewcommand{\contentsname}{\color{gblue}Índice de Contenidos}\tableofcontents}
\clearpage

\setlength{\tabcolsep}{3.5pt}
\newcolumntype{P}[1]{>{\raggedright\arraybackslash}p{#1}}
\newcolumntype{C}[1]{>{\centering\arraybackslash}p{#1}}
\newcommand{\cab}[1]{\textcolor{white}{\textbf{#1}}}

\section{Resumen del estado del proyecto}
A fecha de \hl{22 de octubre de 2026}, el equipo \hl{Los Gantt-sters} ha completado la fase de definición del proyecto \hl{TuAgendita} (Prácticas 0, 1 y 2, propuestas técnica y de solución, manual de coordinación y Product Backlog en Jira) y \hl{arranca hoy el Sprint 1}.

\begin{longtable}{|P{5.2cm}|P{9.8cm}|}
\hline
\rowcolor{glight}\textbf{\color{gblue}Producto} & TuAgendita (proyecto AUTOVIDA, CEE Fundación Purísima Concepción) \\ \hline
\rowcolor{glight}\textbf{\color{gblue}Marco de trabajo} & Scrum con 3 sprints de 3 semanas; gestión en Jira (proyecto SCRUM) \\ \hline
\rowcolor{glight}\textbf{\color{gblue}Tecnología} & Native (front, webapp para tablets), Supabase (backend), Vercel (despliegue), GitHub \\ \hline
\rowcolor{glight}\textbf{\color{gblue}Backlog} & 7 épicas, 26 historias (86 SP), 77 subtareas \\ \hline
\rowcolor{glight}\textbf{\color{gblue}Sprint en curso} & Sprint 1: 22/10/2026 -- 12/11/2026, 12 historias, 34 SP \\ \hline
\rowcolor{glight}\textbf{\color{gblue}Horas registradas hasta hoy} & 59 h de equipo (ver sección~\ref{sec:horas}) \\ \hline
\end{longtable}

\section{Cronograma real de hitos}
\begin{longtable}{|P{2.5cm}|P{9.7cm}|C{2.8cm}|}
\hline
\rowcolor{gblue}\cab{Fecha} & \cab{Hito / actividad} & \cab{Estado} \\ \hline
\endhead
01/10/2026 & Constitución del equipo y primera sesión de trabajo (12 h de equipo registradas). & Completado \\ \hline
08/10/2026 & \textbf{Acta n.º 1.} Reunión de coordinación (8:30--9:00). Se acuerda Supabase + Native + Vercel en lugar de Flutter (A-01) y desarrollar una webapp (A-02). Tareas: acta, propuesta técnica, organización en Jira y propuesta de solución (límite 09/10). & Completado \\ \hline
15/10/2026 & \textbf{Acta n.º 2.} Reunión de seguimiento (8:30--9:00). Se acuerda que los tutores no muevan los botones (A-01) y un único tutor de familia por estudiante (A-02). Tareas de revisión y organización de iteraciones en Jira (límite 16/10). & Completado \\ \hline
16/10/2026 & Fecha límite interna de los documentos, según la norma de finalizar dos días antes de la entrega. & Completado \\ \hline
20--21/10/2026 & Cierre y revisión final de la documentación (Prácticas 0, 1, 2, manual de coordinación, propuestas). & Completado \\ \hline
22/10/2026 & Entrega oficial de la documentación inicial (P0--P2). & Completado \\ \hline
\rowcolor{glight}22/10/2026 & \hl{Inicio del Sprint 1 (Iteración 1):} reunión de arranque y planificación del sprint. & \textbf{En curso} \\ \hline
\end{longtable}

\subsection{Entregables de la fase inicial}
\begin{longtable}{|P{5.7cm}|P{6.3cm}|C{3cm}|}
\hline
\rowcolor{gblue}\cab{Entregable} & \cab{Documento} & \cab{Estado} \\ \hline
\endhead
Práctica 0 & Practica\_0\_Los\_Gantt-sters & Entregado \\ \hline
Práctica 1 (personas y escenarios) & Practica\_1\_Los\_Gantt-sters & Entregado \\ \hline
Práctica 2 (backlog y plan de entregas) & Practica\_2\_Los\_Gantt-sters & Entregado \\ \hline
Manual de coordinación & Manual\_Coordinacion & Entregado \\ \hline
Propuesta técnica y de solución & propuesta\_tecnica, propuesta\_solucion & Entregado \\ \hline
Actas de reunión 1 y 2 & Acta\_08\_10, Acta\_15\_10 & Entregado \\ \hline
Product Backlog en Jira & Proyecto SCRUM (7 épicas, 26 HU, 77 subtareas) & Cargado \\ \hline
\end{longtable}

\section{Registro real de horas}\label{sec:horas}
Datos extraídos del \textit{Registro Horas} del equipo (1 al 21 de octubre de 2026).

\begin{longtable}{|P{3.3cm}|C{1.6cm}|C{1.6cm}|C{1.6cm}|C{1.6cm}|C{1.6cm}|C{1.7cm}|}
\hline
\rowcolor{gblue}\cab{Miembro} & \cab{01/10} & \cab{06/10} & \cab{08/10} & \cab{15/10} & \cab{20--21/10} & \cab{Total} \\ \hline
\endhead
Álvaro Sanabria & 2 & -- & 3 & 3 & 1 + 3 & \textbf{12} \\ \hline
Juan Luis Jiménez & 2 & -- & 3 & 3 & 1 & \textbf{9} \\ \hline
Ana Meca & 2 & -- & 3 & 3 & 1 & \textbf{9} \\ \hline
Jaime Montávez & 2 & -- & 3 & 3 & 1 & \textbf{9} \\ \hline
Germán Delgado & 2 & -- & 3 & 3 & 1 & \textbf{9} \\ \hline
Javier Miñán & 2 & 2 & 3 & 3 & 1 & \textbf{11} \\ \hline
\rowcolor{glight}\textbf{Total equipo} & \textbf{12} & \textbf{2} & \textbf{18} & \textbf{18} & \textbf{9} & \textbf{59} \\ \hline
\end{longtable}

\begin{verde}{Observaciones sobre las horas}
\begin{itemize}
\item Las horas de las reuniones del 08/10 y 15/10 incluyen la sesión de clase y el trabajo conjunto posterior a la reunión de 30 minutos.
\item El 03/10 consta una marca \textit{X} de Jaime Montávez sin horas asociadas.
\item El registro aún no incluye el trabajo no presencial de las semanas previas al sprint: se recomienda que cada miembro lo complete antes del 29/10.
\end{itemize}
\end{verde}

\section{Estimado frente a real}
\begin{longtable}{|P{6cm}|C{4.5cm}|C{4.5cm}|}
\hline
\rowcolor{gblue}\cab{Concepto} & \cab{Estimado} & \cab{Real (hoy)} \\ \hline
\endhead
Dedicación semanal por persona & 7,5 h & 2--4 h (fase documental) \\ \hline
Horas brutas por sprint (6 pers.\ $\times$ 7,5 h $\times$ 3 sem) & 135 h & En curso \\ \hline
Horas ideales por sprint (factor de enfoque 0,65) & 87,75 h & En curso \\ \hline
Velocidad por sprint (2,8 h ideales/SP) & 31 SP & Pendiente de medir \\ \hline
Carga Sprint 1 & 34 SP & 0 SP completados \\ \hline
Horas totales del proyecto (9 semanas) & 405 h & 59 h (14,6\,\%) \\ \hline
\end{longtable}

\section{Planificación real del Sprint 1}
\subsection{Datos del sprint}
\begin{itemize}
\item \textbf{Fechas:} del 22 de octubre al 12 de noviembre de 2026 (3 semanas; las sesiones de clase son los jueves).
\item \textbf{Objetivo:} establecer la \hl{infraestructura técnica} (Supabase, CI/CD, sistema de diseño accesible) y entregar el \hl{núcleo de autonomía personal}: acceso, agenda ``Hecho -- Ahora -- Después'', confirmación de tareas, temporizador, avisos y botón de necesidades.
\item \textbf{Release asociada:} v0.1 (Entrega 1).
\item \textbf{Fin de desarrollo interno:} 10/11/2026, dos días antes de la entrega oficial, según el manual de coordinación.
\end{itemize}

\subsection{Sprint Backlog (12 historias, 34 SP)}
\begin{longtable}{|C{1.7cm}|P{6.6cm}|C{1.5cm}|C{1.7cm}|C{2cm}|}
\hline
\rowcolor{gblue}\cab{ID} & \cab{Historia} & \cab{SP} & \cab{Semana} & \cab{Estado} \\ \hline
\endhead
HUT-01 & Backend y base de datos en Supabase & 3 & 1 & Por hacer \\ \hline
HUT-02 & Pipeline CI/CD y calidad de código & 2 & 1 & Por hacer \\ \hline
HUT-03 & Arquitectura UI y Design System accesible & 3 & 1 & Por hacer \\ \hline
HU-02 & Inicio de sesión estándar para supervisores & 2 & 1 & Por hacer \\ \hline
HU-01 & Inicio de sesión accesible por foto o pictograma & 3 & 2 & Por hacer \\ \hline
HU-03 & Gestión de perfiles y tutores & 3 & 2 & Por hacer \\ \hline
HU-04 & Personalización estética y sensorial (ARASAAC) & 3 & 2 & Por hacer \\ \hline
HU-06 & Agenda ``Hecho -- Ahora -- Después'' & 5 & 2 & Por hacer \\ \hline
HU-08 & Confirmación accesible de inicio y fin de tarea & 2 & 3 & Por hacer \\ \hline
HU-09 & Temporizador visual (Time Timer) & 3 & 3 & Por hacer \\ \hline
HU-16 & Notificación sensorial 5 min antes & 2 & 3 & Por hacer \\ \hline
HU-18 & Botón de necesidades fisiológicas básicas & 3 & 3 & Por hacer \\ \hline
\rowcolor{glight}\multicolumn{2}{|r|}{\textbf{Total}} & \textbf{34} & & \\ \hline
\end{longtable}

\subsection{Calendario semanal previsto}
\begin{longtable}{|P{2.9cm}|P{2.8cm}|P{6.3cm}|C{2.5cm}|}
\hline
\rowcolor{gblue}\cab{Semana} & \cab{Fechas} & \cab{Contenido} & \cab{SP acumulados} \\ \hline
\endhead
Semana 1 & 22/10 -- 29/10 & Infraestructura: HUT-01, HUT-02, HUT-03 y HU-02. Prototipo (\textit{spike}) del canal Realtime de HU-18 (T-18.3). & 10 \\ \hline
Semana 2 & 29/10 -- 05/11 & Perfiles y accesos: HU-01, HU-03, HU-04 y HU-06. & 24 \\ \hline
Semana 3 & 05/11 -- 12/11 & Interacción con la agenda: HU-08, HU-09, HU-16 y HU-18. Congelación del código el 10/11, revisión de calidad y despliegue. & 34 \\ \hline
\end{longtable}

\subsection{Hitos y ceremonias del sprint}
\begin{longtable}{|P{2.5cm}|P{12.5cm}|}
\hline
\rowcolor{gblue}\cab{Fecha} & \cab{Evento} \\ \hline
\endhead
22/10 & Sprint Planning: asignar en Jira las tareas a cada miembro y comprometer los 34 SP. \\ \hline
29/10 & Reunión de seguimiento n.º 3 (revisión de la semana 1, primer despliegue en Vercel). \\ \hline
05/11 & Reunión de seguimiento n.º 4 (revisión de la semana 2, revisión del \textit{burndown}). \\ \hline
10/11 & Congelación del código; revisión del gestor de calidad y auditoría de usabilidad y accesibilidad. \\ \hline
12/11 & Sprint Review (release v0.1), Retrospectiva y Sprint Planning del Sprint 2. \\ \hline
\end{longtable}

\section{Organización del trabajo en el sprint}
\begin{itemize}
\item \textbf{Asignación:} al inicio del sprint se asigna el trabajo de cada persona en Jira. Quien complete menos del 75\,\% de lo asignado acumula la carga pendiente en la iteración siguiente.
\item \textbf{Flujo de ramas:} cada tarea se desarrolla en una rama paralela, el gestor de calidad la revisa y después se fusiona en \texttt{master}.
\item \textbf{Calidad:} ESLint, Prettier, listas de comprobación y registro de no conformidades. El gestor de usabilidad y accesibilidad revisa el código del gestor de calidad.
\item \textbf{Cambios:} cualquier cambio se trata en una reunión específica convocada por el coordinador.
\item \textbf{Roles:} los roles rotan entre reuniones (redactor, moderador, coordinador y presentador).
\end{itemize}

\section{Seguimiento de riesgos}
\begin{longtable}{|P{3.8cm}|C{1.7cm}|C{1.7cm}|P{7.8cm}|}
\hline
\rowcolor{gblue}\cab{Riesgo} & \cab{Prob.} & \cab{Impacto} & \cab{Mitigación} \\ \hline
\endhead
Complejidad de Supabase Realtime (HU-18) & Media & Alto & Prototipo de T-18.3 en la semana 1. \\ \hline
Curva de aprendizaje de Native y Supabase (cambio acordado el 08/10) & Media & Alto & Priorizar HUT-01 y HUT-03; programar en pareja; documentar el \textit{setup}. \\ \hline
Dedicación real inferior a las 7,5 h/semana & Media & Medio & Seguimiento semanal del registro de horas y del \textit{burndown}. \\ \hline
Sprint 1 cargado (34 SP frente a 31 SP de velocidad) & Alta & Medio & Si la velocidad queda por debajo, mover HU-16 o HU-09 al Sprint 2. \\ \hline
Cumplimiento de criterios de accesibilidad (WCAG 2.1 AA) & Baja & Alto & Auditoría T-04.4 y revisión del gestor de usabilidad antes del 10/11. \\ \hline
\end{longtable}

\section{Próximos pasos}
\begin{itemize}
\item Celebrar la reunión de arranque y el Sprint Planning (22/10) y levantar el acta n.º 3.
\item Asignar las subtareas del Sprint 1 en Jira y activar el sprint.
\item Configurar el repositorio de GitHub, el proyecto de Supabase y el proyecto de Vercel.
\item Actualizar el registro de horas con el trabajo no presencial.
\end{itemize}

\section{Conclusiones}
La fase documental se ha cerrado en plazo y con todo el equipo presente en las dos reuniones. El Sprint 1 empieza con 34 SP frente a una velocidad estimada de 31 SP, por lo que el seguimiento semanal del \textit{burndown} y de las horas permitirá reajustar el alcance. Este documento se actualizará al cierre del sprint con la velocidad real.

\end{document}
